% Converted from fixedConstraints.md.
% Required packages: \usepackage{amsmath,longtable,booktabs,array,hyperref}

\section{Constraints (corrected)}\label{constraints-corrected}

This file replaces \texttt{protocol\_constraints.md}. It covers all three phases of the clearing rule: depth accumulation and volume maximization (Phase 1), the imbalance tie-break (Phase 2), and IO admission with pro-rata rationing (Phase 3, Section 7b). The original had no Phase 3 constraints.

Section 11 maps every original constraint (\#1--\#33) to its status in this file. Section 13 records the design decisions taken where the original was silent. Section 12 lists what remains outside the proof.

\textbf{Revision 2} (after the author\textquotesingle s scoping answers): Phase 3 added (F24--F33); the valley pin F20 split by imbalance side; the tie-break F21 changed from "lowest tick" to the rule Phase 3 forces; receipt extended with \(s\), \(P\), \(J\).

\textbf{Revision 3} (as implemented): the no-cross statement is F1--F8 with \(AccA \cdot AccB = 0\) (Section 9).

\textbf{Revision 5} (as implemented): the update-validity proof (Section 9b), a second proof per batch showing that every logged update adds a non-negative quantity at one tick.

\textbf{Revision 4} (as implemented): range checks use limbs and a lookup argument in place of one committed column per bit (Section 2). F23 changes from booleanity to table membership. No other constraint changes, and the ranges proved are the same. The bit-column text is kept in \texttt{archive/fixedConstraints\_rev3\_bit-columns.md}.

\begin{center}\rule{0.5\linewidth}{0.5pt}\end{center}

\subsection{Setting}\label{setting}

\textbf{Domain.} \(H = \{1, \omega, \ldots, \omega^{n-1}\}\), with \(\omega\) a primitive \(n\)-th root of unity and \(n\) a power of two. The book has \(N \le n\) real price ticks; tick \(i\) sits at \(\omega^i\). Ticks \(N, \ldots, n-1\) are padding. \(Z_H(X) = X^n - 1\).

\textbf{Vanishing equations.} We write \(V(X) = 0\) to mean \(Z_H(X) \mid V(X)\). The prover commits to the quotient; the verifier checks the identity at a random \(\zeta \notin H\).

\textbf{Lagrange basis.} \(L_i(X)\) is 1 at \(\omega^i\) and 0 on the rest of \(H\). All indices run from 0 to \(n-1\). We use \(L_i\) for every single-point gate, including the accumulator initializations.

\textbf{Bit widths.} Two widths are fixed at setup and do not depend on any disclosed value:

\begin{itemize}
\item
  \(k_{in}\): width of a single tick\textquotesingle s volume. Every \(B_i, A_i \in [0, 2^{k_{in}})\).
\item
  \(k\): width of every accumulated quantity, with \(k \ge k_{in} + \lceil \log_2 N \rceil\). This guarantees that total volume on either side is below \(2^k\), so every honest surplus, gap, and slack fits in \(k\) bits. We also require \(2^{k+1} < q\), which holds trivially for BN254 at any practical \(k\).
\end{itemize}

\textbf{Limbs.} Range-checked quantities are written in limbs of \(w\) bits, with \(w = \min(8, \log_2 n)\), so that the table \(\{0, \ldots, 2^w - 1\}\) fits in the domain. A single tick\textquotesingle s volume has \(m_{in} = \lceil k_{in}/w \rceil\) limbs and an accumulated quantity has \(m = \lceil k/w \rceil\).

\textbf{Public scalars (the clearing receipt).}

\begin{longtable}[]{@{}ll@{}}
\toprule\noalign{}
Scalar & Meaning \\
\midrule\noalign{}
\endhead
\bottomrule\noalign{}
\endlastfoot
\(V_{max}\) & Maximum executable volume \\
\(c, d\) & First and last tick of the plateau \\
\(p^*\) & Clearing tick \\
\(V_{min\Delta}\) & Imbalance at \(p^*\) (minimum over the plateau) \\
\(s \in \{0,1\}\) & Imbalance side: \(0\) = demand long (buy imbalance), \(1\) = supply long. Set to \(0\) when \(V_{min\Delta} = 0\) \\
\(P\) & Volume resting exactly at \(p^*\) on the long side. Set to \(0\) when \(V_{min\Delta} = 0\) \\
\(J\) & IO volume on the short side eligible at \(p^*\). Set to \(0\) when \(V_{min\Delta} = 0\) \\
\end{longtable}

\textbf{Verifier-side scalar checks.} Before any polynomial check, the verifier rejects unless

\[0 \le c \le p^* \le d \le N-1, \qquad 1 \le V_{max} < 2^k, \qquad 0 \le V_{min\Delta} < 2^k,\]

\[s \in \{0,1\}, \qquad V_{min\Delta} \le P < 2^k, \qquad 0 \le J < 2^k, \qquad V_{min\Delta} = 0 \Rightarrow s = P = J = 0.\]

The condition \(c \le p^* \le d\) is load-bearing: no polynomial constraint forces the clearing tick into the plateau. The condition \(V_{min\Delta} \le P\) is load-bearing for the tie-break (Section 7). The case \(V_{max} = 0\) (no cross) is handled separately in Section 9.

\textbf{Derived by the verifier} (no proof needed, public arithmetic):

\[I = \min(J, V_{min\Delta}), \qquad R = V_{min\Delta} - I, \qquad V_{exec} = V_{max} + I.\]

\(I\) is the admitted IO volume, \(R\) the residual imbalance left unfilled on the long side, \(V_{exec}\) the total traded volume.

We write \(\nu = [V_{min\Delta} > 0]\) for the public bit that says an imbalance exists.

\textbf{Public polynomials.} The verifier computes these from the receipt; none is committed.

\begin{longtable}[]{@{}lll@{}}
\toprule\noalign{}
Polynomial & Definition & Role \\
\midrule\noalign{}
\endhead
\bottomrule\noalign{}
\endlastfoot
\(Pad(X)\) & \(\sum_{i=N}^{n-1} L_i(X)\) (zero if \(N = n\)) & Marks padding ticks \\
\(Mask_P(X)\) & \(\sum_{i=c}^{d} L_i(X)\) & Marks the plateau \\
\(Mask_C(X)\) & \([c \ge 1]\,L_{c-1}(X) + [d \le n-2]\,L_{d+1}(X)\) & Marks the cliff ticks that exist \\
\(Mask_T(X)\) & \([p^* > c]\cdot[\,s = 1 \text{ or } V_{min\Delta} = 0\,]\cdot L_{p^*-1}(X)\) & Marks the tick just below \(p^*\), when the tie-break needs it \\
\(T(X)\) & \(\sum_{i=0}^{2^w-1} i \cdot L_i(X)\) & The lookup table: value \(i\) at \(\omega^i\) for \(i < 2^w\), zero elsewhere. Depends on \(n\) only \\
\end{longtable}

The bracket \([\cdot]\) is 1 if the condition holds and 0 otherwise. The verifier evaluates it on public scalars, so a cliff that does not exist is dropped by the verifier, not by an index that wraps around the domain.

\textbf{Committed columns.}

\begin{longtable}[]{@{}lll@{}}
\toprule\noalign{}
Column & Polynomial & Meaning \\
\midrule\noalign{}
\endhead
\bottomrule\noalign{}
\endlastfoot
Bids++ & \(B(X)\) & Raw bid volume per tick \\
Asks++ & \(A(X)\) & Raw ask volume per tick \\
Bid Depth & \(AccB(X)\) & Cumulative demand, summed from the top tick down \\
Ask Depth & \(AccA(X)\) & Cumulative supply, summed from the bottom tick up \\
Min(Bid,Ask) & \(Min(X)\) & Executable volume \\
Bid Surplus++ & \(SurpB(X)\) & \(AccB - Min\) \\
Ask Surplus++ & \(SurpA(X)\) & \(AccA - Min\) \\
Abs(Delta) & \(\Delta(X)\) & \(SurpA + SurpB\) \\
Slack++ & \(Slack(X)\) & One column serving both cliffs and the tie-break \\
IO Bids & \(IOB(X)\) & Raw imbalance-only bid volume per tick \\
IO Asks & \(IOA(X)\) & Raw imbalance-only ask volume per tick \\
IO Bid Depth & \(IOAccB(X)\) & IO bids with limit at or above the tick \\
IO Ask Depth & \(IOAccA(X)\) & IO asks with limit at or below the tick \\
Limb columns & \(\ell^{B}_l, \ell^{A}_l, \ell^{IOB}_l, \ell^{IOA}_l\) (\(m_{in}\) each); \(\ell^{sB}_l, \ell^{sA}_l, \ell^{ceil}_l, \ell^{flr}_l, \ell^{slk}_l\) (\(m\) each) & Range witnesses \\
Multiplicities & \(M(X)\) & \(M(\omega^v)\) = number of looked-up cells equal to \(v\) \\
Lookup helpers & \(h_1, \ldots, h_{\lceil E/2 \rceil}\), \(U(X)\), \(S(X)\) & Committed after the challenge \(\beta\) (Section 2) \\
\end{longtable}

The first prover round has 13 value columns, \(4m_{in} + 5m\) limb columns and \(M\). The second round has \(\lceil E/2 \rceil + 2\) columns, where \(E\) is the number of looked-up values: \(E = 4m_{in} + 5m\), plus 4 if \(w \nmid k_{in}\), plus 5 if \(w \nmid k\). For \(n = 1024\), \(k_{in} = 32\), \(k = 42\): \(w = 8\), 46 limb columns, \(E = 51\), 26 helpers, 88 columns in all (the bit-column version had 351).

\begin{center}\rule{0.5\linewidth}{0.5pt}\end{center}

\subsection{The Range Gadget}\label{the-range-gadget}

We use one gadget for every inequality. To prove \(P(\omega^i) \in [0, 2^{b})\) at every tick, with \(m_P = \lceil b/w \rceil\), the prover supplies limb columns \(\ell_0, \ldots, \ell_{m_P-1}\) with

\[P(X) - \sum_{l=0}^{m_P-1} 2^{wl}\, \ell_l(X) = 0,\]

and proves that every limb value lies in \(\{0, \ldots, 2^w-1\}\). If \(b\) is not a multiple of \(w\), write \(r = b \bmod w\); the top limb must then be below \(2^r\), which holds exactly when \(2^{w-r}\,\ell_{m_P-1}\) is also in the table. So the \emph{looked-up values} of \(P\) are its \(m_P\) limbs and, when \(r \ne 0\), the scaled top limb. The range proved is exactly \([0, 2^b)\), as with bit columns.

A value that is negative as an integer appears in the field as an element near \(q\), which has no such representation when \(2^{b} < q\).

\textbf{Table membership (logUp).} Let \(f_1, \ldots, f_E\) be all looked-up values of all range-checked quantities; each is a limb column times a public power of two. The prover commits to \(M\) together with the other first-round columns. The verifier\textquotesingle s challenge \(\beta\) is drawn after those commitments. The prover then commits to the helpers, defined on \(H\) by

\[h_j = \frac{1}{\beta - f_{2j-1}} + \frac{1}{\beta - f_{2j}}, \qquad U = \frac{M}{\beta - T}, \qquad S(\omega^{i+1}) = S(\omega^{i}) + \sum_j h_j(\omega^i) - U(\omega^i), \quad S(1) = 0.\]

The identities, each on all of \(H\), are

\[h_j(X)\,\big(\beta - f_{2j-1}(X)\big)\big(\beta - f_{2j}(X)\big) - \big(\beta - f_{2j-1}(X)\big) - \big(\beta - f_{2j}(X)\big) = 0 \quad \text{for each pair},\]

\[U(X)\,\big(\beta - T(X)\big) - M(X) = 0,\]

\[S(\omega X) - S(X) - \Big(\sum_j h_j(X) - U(X)\Big) = 0.\]

If \(E\) is odd the last helper covers one value: \(h(X)(\beta - f_E(X)) - 1 = 0\).

The last identity has no gate at \(\omega^{n-1}\): it also holds for the step from the last point back to the first. Summing it over \(H\), the left side telescopes to zero around the cycle, which gives

\[\sum_{i=0}^{n-1} \sum_{e=1}^{E} \frac{1}{\beta - f_e(\omega^i)} = \sum_{i=0}^{n-1} \frac{M(\omega^i)}{\beta - T(\omega^i)}.\]

No initial condition on \(S\) is needed; only its differences are constrained, and \(S(1)\) may be any value.

\textbf{Soundness of the lookup.} We use the following fact (Haböck, ePrint 2022/1530, Lemma 5): over a field of characteristic \(p > N\), for sequences \((a_i)_{i<N}\) and \((b_i)_{i<N}\), \(\{a_i\} \subseteq \{b_i\}\) if and only if there are field elements \((m_i)\) with \(\sum_i 1/(X - a_i) = \sum_i m_i/(X - b_i)\) as rational functions. The multiplicities are arbitrary field elements; nothing constrains them to be counts. The direction we need follows from uniqueness of partial fractions: comparing the residue at a value \(z\), the number of \(i\) with \(a_i = z\) equals \(\sum_{i: b_i = z} m_i\) in the field. If \(z\) is not a table value the right side is empty, so the count must vanish in the field; the count is a positive integer no larger than the number of looked-up cells, which is \(En < q\), so it cannot. Here the \(a_i\) are the \(En\) cells \(f_e(\omega^i)\) and the \(b_i\) are the \(n\) values of \(T\) on \(H\) (the value 0 repeats on the points beyond \(\omega^{2^w-1}\); repeats are allowed).

The limbs and \(M\) are committed before \(\beta\) is drawn. Suppose some cell is not a table value. Then \(\sum_{i,e} 1/(X - f_e(\omega^i)) - \sum_i M(\omega^i)/(X - T(\omega^i))\) is a non-zero rational function whose numerator has degree below \((E+1)n\). A cheating prover can satisfy the three identities on \(H\) only if \(\beta\) falls in one of these sets:

\begin{enumerate}
\def\labelenumi{\arabic{enumi}.}
\item
  the roots of that numerator: fewer than \((E+1)n\) values;
\item
  the cell values: at most \(En\) values. If \(\beta\) equals both looked-up values of a pair at some point of \(H\), the pair identity reads \(0 = 0\) there and leaves \(h_j\) free at that point;
\item
  the table values: at most \(2^w \le n\) values. If \(\beta = T(\omega^i)\) and \(M(\omega^i) = 0\), the table-helper identity leaves \(U(\omega^i)\) free.
\end{enumerate}

Outside these sets every denominator is non-zero, the identities force \(h_j\) and \(U\) to be the stated fractions on \(H\), and the telescoped sum is the non-zero rational function evaluated at \(\beta\). So the lookup adds at most \((2E+2)\,n/q\) to the soundness error; for \(E = 51\) and \(n = 2^{13}\) this is below \(2^{-230}\). The errors of the other steps (the combination challenge \(\alpha\), the evaluation point \(\zeta\), the opening batching) are the usual ones and are unchanged.

An earlier revision of this paragraph gave the bound \((E+1)n\) and omitted cases 2 and 3.

\textbf{Completeness.} The honest prover fails only if \(\beta\) equals a cell value or a table value, since a denominator is then zero: probability at most \((E+1)\,n/q\). The implementation does not resample; the proof would fail to be produced (\texttt{NotDivisible}).

\textbf{Sources checked (2026-11-18).} The ePrint server refused automated access, so the paper\textquotesingle s PDF was not read directly. The statement of Lemma 5 above (hypothesis \(p > N\), multiplicities as field elements, the equivalence) was confirmed from an indexed excerpt of the paper and from two independent restatements: the ArkLib formalization of Haböck\textquotesingle s Lemmas 3 and 5 (which notes that the bound that matters is on the largest count of a single value, weaker than what we assume), and "Formal verification of the S-two AIR" (arXiv 2606.04311), which uses the same lemma, the same cyclic running sum with a free starting value, and the error bound (number of fractions + rows)\(/|F|\) for the challenge. The ordering "commit multiplicities, then draw the challenge, then commit helpers" is the one used in cq and in the multi-column logUp protocol. The author should read Section 3 of the paper once before the security section is final and cite the lemma number from the version cited.
Points of \(H\) beyond \(\omega^{2^w-1}\) carry the table value 0, which is in the table, so they add nothing.

\textbf{Cost.} One lookup, shared by all nine range-checked quantities: one multiplicity column, one helper per two looked-up values, \(U\) and \(S\). The pair identity has degree 3 in committed columns; every other constraint has degree at most 2. \(S\) is opened at \(\zeta\) and \(\omega\zeta\).

\begin{center}\rule{0.5\linewidth}{0.5pt}\end{center}

\subsection{\texorpdfstring{Inputs: \(B(X)\), \(A(X)\)}{Inputs: B(X), A(X)}}\label{inputs-bx-ax}

\textbf{F1, F2 --- Input range.}

\[B(X) - \sum_{l=0}^{m_{in}-1} 2^{wl}\, \ell^{B}_l(X) = 0, \qquad A(X) - \sum_{l=0}^{m_{in}-1} 2^{wl}\, \ell^{A}_l(X) = 0\]

\textbf{F3, F4 --- Padding is empty.}

\[Pad(X) \cdot B(X) = 0, \qquad Pad(X) \cdot A(X) = 0\]

Without F3, volume placed on a padding tick \(i \ge N\) acts as a bid above every real price, which is a hidden market buy. F4 is included for symmetry: ask volume on a padding tick affects no real tick, but we do not want the statement to depend on that observation.

Market orders are recorded at the boundary real ticks: a market buy at tick \(N-1\), a market sell at tick \(0\).

\begin{center}\rule{0.5\linewidth}{0.5pt}\end{center}

\subsection{\texorpdfstring{Depth: \(AccB(X)\), \(AccA(X)\)}{Depth: AccB(X), AccA(X)}}\label{depth-accbx-accax}

\textbf{F5 --- Demand initialization.} \(L_{n-1}(X) \cdot \big(AccB(X) - B(X)\big) = 0\)

\textbf{F6 --- Supply initialization.} \(L_{0}(X) \cdot \big(AccA(X) - A(X)\big) = 0\)

\textbf{F7 --- Demand transition.} \(AccB[i] = AccB[i+1] + B[i]\):

\[(X - \omega^{n-1}) \cdot \big[AccB(X) - AccB(\omega X) - B(X)\big] = 0\]

\textbf{F8 --- Supply transition.} \(AccA[i+1] = AccA[i] + A[i+1]\):

\[(X - \omega^{n-1}) \cdot \big[AccA(\omega X) - AccA(X) - A(\omega X)\big] = 0\]

\textbf{Correction.} The original \#6 used \(A(X)\) in place of \(A(\omega X)\). With the initialization \(AccA[0] = A[0]\), that version counts \(A[0]\) twice and never counts \(A[n-1]\). The two recurrences are not mirror images in their arguments: the demand recurrence adds the volume at the current tick, the supply recurrence adds the volume at the next tick.

\textbf{Padding.} Both recurrences run over the whole domain and are gated only at the wraparound point \(\omega^{n-1}\). With F3 and F4, \(AccA\) is constant across the padding region (equal to total supply) and \(AccB\) is zero there. No second gating point is needed. A forward accumulator that appears to break at tick \(N-1\) is the symptom of zero-filling the derived column \(AccA\) instead of extending the running sum through the padding.

\textbf{Monotonicity.} F1 and F7 force \(AccB\) to be non-increasing in the tick index; F2 and F8 force \(AccA\) to be non-decreasing. Sections 6 and 7 rely on both facts.

\begin{center}\rule{0.5\linewidth}{0.5pt}\end{center}

\subsection{\texorpdfstring{Executable Volume: \(Min(X)\), \(SurpB(X)\), \(SurpA(X)\)}{Executable Volume: Min(X), SurpB(X), SurpA(X)}}\label{executable-volume-minx-surpbx-surpax}

\textbf{F9, F10 --- Surplus definitions.}

\[SurpB(X) - \big(AccB(X) - Min(X)\big) = 0, \qquad SurpA(X) - \big(AccA(X) - Min(X)\big) = 0\]

\textbf{Correction.} The original list bit-decomposed \(SurpB\) and \(SurpA\) but never tied them to \(AccB - Min\) and \(AccA - Min\). Without F9 and F10 the prover may set both surpluses to zero everywhere, and the non-negativity checks prove nothing about \(Min\).

\textbf{F11 --- Min is one of the two depths.}

\[SurpA(X) \cdot SurpB(X) = 0\]

This is the original \#7, \((AccA - Min)(AccB - Min) = 0\), rewritten through F9 and F10.

\textbf{F12, F13 --- Surplus non-negativity.}

\[SurpB(X) - \sum_{l=0}^{m-1} 2^{wl}\, \ell^{sB}_l(X) = 0, \qquad SurpA(X) - \sum_{l=0}^{m-1} 2^{wl}\, \ell^{sA}_l(X) = 0\]

F9--F13 together force \(Min(\omega^i) = \min(AccA[i], AccB[i])\) at every tick: F11 puts \(Min\) on one of the two curves, and F12, F13 rule out the larger one.

\textbf{F14 --- Ceiling.}

\[\big(V_{max} - Min(X)\big) - \sum_{l=0}^{m-1} 2^{wl}\, \ell^{ceil}_l(X) = 0\]

\begin{center}\rule{0.5\linewidth}{0.5pt}\end{center}

\subsection{\texorpdfstring{Phase 1: the Plateau \([c, d]\)}{Phase 1: the Plateau {[}c, d{]}}}\label{phase-1-the-plateau-c-d}

\textbf{F15, F16 --- Endpoint pins.}

\[L_c(X) \cdot \big(Min(X) - V_{max}\big) = 0, \qquad L_d(X) \cdot \big(Min(X) - V_{max}\big) = 0\]

\textbf{F17 --- Cliffs.} At each cliff tick that exists, \(Min\) is strictly below \(V_{max}\):

\[Mask_C(X) \cdot \big(V_{max} - Min(X) - 1 - Slack(X)\big) = 0\]

\(Slack\) is range-checked by F22.

\textbf{Corrections.}

\begin{itemize}
\item
  The original \(Mask_C\) was constrained (booleanity, disjointness from \(Mask_P\)) but never used; the cliff equations were gated by \(L_{c-1}\) and \(L_{d+1}\) directly. We gate with the public \(Mask_C\) and drop the unused constraints.
\item
  When \(c = 0\) the index \(c-1\) wraps to \(n-1\), and when \(d = n-1\) the index \(d+1\) wraps to \(0\). The original equations then pin an unrelated tick. \(Mask_C\) omits the missing cliff.
\item
  When \(d = N-1\) and \(N < n\), tick \(d+1\) is a padding tick where \(Min = 0\), so the right cliff holds for every honest book with \(V_{max} \ge 1\).
\item
  The original both opened \(Min\) at the cliff ticks in plaintext and hid the gap behind a slack. We keep the slack and remove the openings. Revealing \(Min(\omega^{c-1})\) and \(Min(\omega^{d+1})\) discloses depth information the receipt does not otherwise contain.
\item
  One \(Slack\) column replaces \(Slack_L\) and \(Slack_R\). The cliff ticks are distinct points, so one polynomial carries both values and the range-witness columns are halved.
\end{itemize}

\textbf{Why this pins the plateau.} \(Min\) is the pointwise minimum of a non-increasing and a non-decreasing sequence, so for \(i < j < l\) we have \(Min[j] \ge \min(Min[i], Min[l])\). F14 gives \(Min \le V_{max}\) everywhere. F15 and F16 then give \(Min = V_{max}\) on all of \([c,d]\). If some \(i < c-1\) had \(Min[i] = V_{max}\), then \(Min[c-1] \ge \min(Min[i], Min[c]) = V_{max}\), contradicting F17; the right side is symmetric. So \(V_{max} = \max_i Min[i]\) and \([c,d]\) is exactly the set of maximizing ticks.

\begin{center}\rule{0.5\linewidth}{0.5pt}\end{center}

\subsection{\texorpdfstring{Phase 2: the Clearing Tick \(p^*\)}{Phase 2: the Clearing Tick p\^{}*}}\label{phase-2-the-clearing-tick-p}

\textbf{F18 --- Imbalance.}

\[\Delta(X) - \big(SurpA(X) + SurpB(X)\big) = 0\]

By F11--F13 at least one surplus is zero and both are non-negative, so \(\Delta[i] = |AccA[i] - AccB[i]|\).

\textbf{F19 --- Floor.} No plateau tick has imbalance below \(V_{min\Delta}\):

\[Mask_P(X) \cdot \Big[\big(\Delta(X) - V_{min\Delta}\big) - \sum_{l=0}^{m-1} 2^{wl}\, \ell^{flr}_l(X)\Big] = 0\]

\textbf{F20a, F20b --- Valley pin, by side.} The floor is attained at \(p^*\), on the disclosed side:

\[L_{p^*}(X) \cdot \big(SurpB(X) - (1-s)\,V_{min\Delta}\big) = 0, \qquad L_{p^*}(X) \cdot \big(SurpA(X) - s\,V_{min\Delta}\big) = 0\]

With F18 these give \(\Delta[p^*] = V_{min\Delta}\) and fix which side is long. They replace the single pin \(L_{p^*}(\Delta - V_{min\Delta})\) of revision 1.

\textbf{The tie-break rule.} Phase 3 states that orders strictly better than \(p^*\) execute in full and only at-price orders on the long side are rationed. This is satisfiable only if the whole imbalance sits at the clearing tick: \(V_{min\Delta} \le P\), where \(P\) is the long side\textquotesingle s volume at exactly \(p^*\). We take this as the tertiary rule:

\begin{quote}
Among the plateau ticks with minimum imbalance, \(p^*\) is the lowest one whose at-price long-side volume covers the imbalance.
\end{quote}

\emph{Why a plain "lowest tick" rule is wrong.} Take bids of 10 at tick 5 and asks of 6 at tick 2. Every tick from 2 to 5 executes 6 with imbalance 4. At \(p^* = 2\) the bids at tick 5 are strictly better than the price, yet only 6 of 10 can fill. At \(p^* = 5\) those bids are at-price and are rationed 6/10, and every ask fills. The rule selects 5.

\emph{Candidates.} Let \(g[i] = AccA[i] - AccB[i]\), non-decreasing in \(i\), and \(v = V_{min\Delta}\).

\begin{itemize}
\item
  If \(v = 0\), every minimizing tick satisfies the rule; the lowest one is selected.
\item
  If \(v > 0\), at most two ticks satisfy it: the highest demand-long minimizer (the last tick with \(g = -v\)) and the lowest supply-long minimizer (the first tick with \(g = +v\)). For a demand-long minimizer \(j\) that is not the highest, \(g[j+1] = g[j]\) forces \(B[j] = 0 < v\). The supply side is symmetric. The demand-long candidate, when it exists, is the lower of the two.
\item
  A candidate always exists: if the highest demand-long minimizer \(j\) had \(AccB[j+1] > AccA[j]\), tick \(j+1\) would be a demand-long minimizer too.
\end{itemize}

\emph{Enforcement.} The scalar check \(V_{min\Delta} \le P\) together with F20a/b and F32 (Section 7b) makes \(p^*\) a candidate on side \(s\). It remains to exclude a lower candidate:

\begin{itemize}
\item
  \(s = 0\), \(v > 0\): none can exist. No further check.
\item
  \(s = 1\) or \(v = 0\): a lower candidate exists exactly when \(p^* - 1\) is a plateau tick with imbalance \(v\). F21 excludes it.
\end{itemize}

\textbf{F21 --- Tie-break.}

\[Mask_T(X) \cdot \big(\Delta(X) - V_{min\Delta} - 1 - Slack(X)\big) = 0\]

\(Mask_T\) is non-zero only when \(p^* > c\) and (\(s = 1\) or \(V_{min\Delta} = 0\)). In the demand-long case the tick below \(p^*\) may legitimately have the same imbalance, so the check is switched off.

The tick \(p^*-1\) lies inside the plateau and the cliff ticks lie outside it, so \(Slack\) serves F17 and F21 without collision.

\textbf{Correction.} The original admitted any \(p^*\) that attains \(V_{min\Delta}\). When several plateau ticks tie, the exchange could choose among them.

\textbf{F22 --- Slack non-negativity.}

\[Slack(X) - \sum_{l=0}^{m-1} 2^{wl}\, \ell^{slk}_l(X) = 0\]

The prover sets \(Slack = 0\) at every tick that neither F17 nor F21 touches.

\textbf{F23 --- Table membership.} Every looked-up value of F1, F2, F12, F13, F14, F19, F22, F24, F25 lies in \(\{0, \ldots, 2^w-1\}\). This is the lookup of Section 2: its pair identities, the table-helper identity and the running-sum identity. It replaces the booleanity constraint of revision 3.

\begin{center}\rule{0.5\linewidth}{0.5pt}\end{center}

\subsection{Phase 3: IO Admission and Rationing}\label{phase-3-io-admission-and-rationing}

\textbf{The rule.} Let \(v = V_{min\Delta}\). Take the demand-long case (\(s = 0\)); the supply-long case swaps the roles of bids and asks.

\begin{enumerate}
\def\labelenumi{\arabic{enumi}.}
\item
  Every ask at or below \(p^*\) executes in full (the short side), total \(V_{max}\).
\item
  Every bid strictly above \(p^*\) executes in full.
\item
  Sell IO orders with limit at or below \(p^*\) are eligible, total \(J\). The exchange admits \(I = \min(J, v)\) of them. If \(J > v\), eligible IO orders are filled pro-rata at \(I / J\).
\item
  The residual \(R = v - I\) is left unfilled among the bids at exactly \(p^*\), total \(P\). They are filled pro-rata at \((P - R)/P\).
\item
  Total traded volume is \(V_{exec} = V_{max} + I\).
\end{enumerate}

IO orders never affect \(p^*\): Phases 1 and 2 use limit and market orders only.

\textbf{F24, F25 --- IO input range.}

\[IOB(X) - \sum_{l=0}^{m_{in}-1} 2^{wl}\, \ell^{IOB}_l(X) = 0, \qquad IOA(X) - \sum_{l=0}^{m_{in}-1} 2^{wl}\, \ell^{IOA}_l(X) = 0\]

\textbf{F26, F27 --- IO padding.} \(Pad(X) \cdot IOB(X) = 0\), \(\;Pad(X) \cdot IOA(X) = 0\).

\textbf{F28, F29 --- IO bid depth} (suffix sum, as F5, F7).

\[L_{n-1}(X)\big(IOAccB(X) - IOB(X)\big) = 0, \qquad (X - \omega^{n-1})\big[IOAccB(X) - IOAccB(\omega X) - IOB(X)\big] = 0\]

\textbf{F30, F31 --- IO ask depth} (prefix sum, as F6, F8).

\[L_{0}(X)\big(IOAccA(X) - IOA(X)\big) = 0, \qquad (X - \omega^{n-1})\big[IOAccA(\omega X) - IOAccA(X) - IOA(\omega X)\big] = 0\]

\textbf{F32 --- At-price long volume.}

\[\nu \cdot L_{p^*}(X) \cdot \big((1-s)\,B(X) + s\,A(X) - P\big) = 0\]

\textbf{F33 --- Eligible IO volume.}

\[\nu \cdot L_{p^*}(X) \cdot \big((1-s)\,IOAccA(X) + s\,IOAccB(X) - J\big) = 0\]

When demand is long the short side is supply, so the eligible IO orders are sell IOs with limit at or below \(p^*\), which is \(IOAccA[p^*]\); the supply-long case reads \(IOAccB[p^*]\). Both pins are switched off by \(\nu\) when there is no imbalance, so a balanced batch discloses neither \(P\) nor \(J\).

\textbf{What a trader derives.} From the public receipt alone, a trader computes their own fill:

\begin{longtable}[]{@{}ll@{}}
\toprule\noalign{}
Order & Fill \\
\midrule\noalign{}
\endhead
\bottomrule\noalign{}
\endlastfoot
Limit or market order strictly better than \(p^*\), or on the short side at \(p^*\) & Full \\
Limit order at exactly \(p^*\) on the long side & \(q \cdot (P - R)/P\) \\
IO order on the short side, eligible at \(p^*\) & \(q \cdot I/J\) \\
Anything else & None \\
\end{longtable}

\textbf{Disclosure.} \(s\), \(P\) and \(J\) are new public values. A rationed trader learns the fill ratio from their own execution, and with it \(P\) or \(J\); publishing them gives every participant what the rationed ones already see. Real closing auctions publish the imbalance side and quantity.

\begin{center}\rule{0.5\linewidth}{0.5pt}\end{center}

\subsection{Full Constraint List}\label{full-constraint-list}

\begin{longtable}[]{@{}llll@{}}
\toprule\noalign{}
\# & Name & Equation & Bites at \\
\midrule\noalign{}
\endhead
\bottomrule\noalign{}
\endlastfoot
F1 & Bid range & \(B - \sum 2^{wl} \ell^{B}_l = 0\) & All ticks \\
F2 & Ask range & \(A - \sum 2^{wl} \ell^{A}_l = 0\) & All ticks \\
F3 & Bid padding & \(Pad \cdot B = 0\) & Padding ticks \\
F4 & Ask padding & \(Pad \cdot A = 0\) & Padding ticks \\
F5 & Demand init & \(L_{n-1} \cdot (AccB - B) = 0\) & \(\omega^{n-1}\) \\
F6 & Supply init & \(L_0 \cdot (AccA - A) = 0\) & \(\omega^0\) \\
F7 & Demand transition & \((X-\omega^{n-1})[AccB(X) - AccB(\omega X) - B(X)] = 0\) & All but last \\
F8 & Supply transition & \((X-\omega^{n-1})[AccA(\omega X) - AccA(X) - A(\omega X)] = 0\) & All but last \\
F9 & \(SurpB\) definition & \(SurpB - (AccB - Min) = 0\) & All ticks \\
F10 & \(SurpA\) definition & \(SurpA - (AccA - Min) = 0\) & All ticks \\
F11 & Min selection & \(SurpA \cdot SurpB = 0\) & All ticks \\
F12 & \(SurpB \ge 0\) & \(SurpB - \sum 2^{wl} \ell^{sB}_l = 0\) & All ticks \\
F13 & \(SurpA \ge 0\) & \(SurpA - \sum 2^{wl} \ell^{sA}_l = 0\) & All ticks \\
F14 & Ceiling & \((V_{max} - Min) - \sum 2^{wl} \ell^{ceil}_l = 0\) & All ticks \\
F15 & Plateau left endpoint & \(L_c \cdot (Min - V_{max}) = 0\) & \(\omega^c\) \\
F16 & Plateau right endpoint & \(L_d \cdot (Min - V_{max}) = 0\) & \(\omega^d\) \\
F17 & Cliffs & \(Mask_C \cdot (V_{max} - Min - 1 - Slack) = 0\) & \(\omega^{c-1}\), \(\omega^{d+1}\) where they exist \\
F18 & \(\Delta\) definition & \(\Delta - (SurpA + SurpB) = 0\) & All ticks \\
F19 & Delta floor & \(Mask_P \cdot [(\Delta - V_{min\Delta}) - \sum 2^{wl} \ell^{flr}_l] = 0\) & Plateau \\
F20a & Valley pin, demand side & \(L_{p^*} \cdot (SurpB - (1-s)V_{min\Delta}) = 0\) & \(\omega^{p^*}\) \\
F20b & Valley pin, supply side & \(L_{p^*} \cdot (SurpA - s\,V_{min\Delta}) = 0\) & \(\omega^{p^*}\) \\
F21 & Tie-break & \(Mask_T \cdot (\Delta - V_{min\Delta} - 1 - Slack) = 0\) & \(\omega^{p^*-1}\), when \(p^* > c\) and (\(s=1\) or \(V_{min\Delta}=0\)) \\
F22 & \(Slack \ge 0\) & \(Slack - \sum 2^{wl} \ell^{slk}_l = 0\) & All ticks \\
F23 & Table membership & Lookup identities of Section 2 over all limbs & All ticks \\
F24 & IO bid range & \(IOB - \sum 2^{wl} \ell^{IOB}_l = 0\) & All ticks \\
F25 & IO ask range & \(IOA - \sum 2^{wl} \ell^{IOA}_l = 0\) & All ticks \\
F26 & IO bid padding & \(Pad \cdot IOB = 0\) & Padding ticks \\
F27 & IO ask padding & \(Pad \cdot IOA = 0\) & Padding ticks \\
F28 & IO demand init & \(L_{n-1} \cdot (IOAccB - IOB) = 0\) & \(\omega^{n-1}\) \\
F29 & IO demand transition & \((X-\omega^{n-1})[IOAccB(X) - IOAccB(\omega X) - IOB(X)] = 0\) & All but last \\
F30 & IO supply init & \(L_0 \cdot (IOAccA - IOA) = 0\) & \(\omega^0\) \\
F31 & IO supply transition & \((X-\omega^{n-1})[IOAccA(\omega X) - IOAccA(X) - IOA(\omega X)] = 0\) & All but last \\
F32 & At-price long volume & \(\nu\, L_{p^*} \cdot ((1-s)B + sA - P) = 0\) & \(\omega^{p^*}\), when \(V_{min\Delta} > 0\) \\
F33 & Eligible IO volume & \(\nu\, L_{p^*} \cdot ((1-s)IOAccA + s\,IOAccB - J) = 0\) & \(\omega^{p^*}\), when \(V_{min\Delta} > 0\) \\
\end{longtable}

Plus the verifier\textquotesingle s scalar checks from Section 1. The count is 34 constraints (F20 is two). F23 is one constraint made of \(\lceil E/2 \rceil + 2\) polynomial identities; all identities are combined into one quotient with powers of a challenge \(\alpha\).

\textbf{Rounds.} The prover commits to the value columns, the limb columns and \(M\); receives \(\beta\); commits to the lookup helpers; receives \(\alpha\); commits to the quotient; receives \(\zeta\).

\textbf{Degrees.} Before blinding, every committed column has degree at most \(n-1\). The pair identities of F23 have degree at most \(3n-3\); every other \(V_i\) has degree at most \(2n-2\) (the four accumulator transitions have degree \(n\)). The quotient has degree at most \(2n-3\) before blinding and at most \(2n+3\) with it, so the SRS must reach degree \(2n+3\); the implementation asks for \(2n+8\). Blinding is described in Section 10.

\textbf{Openings.} \(AccB\), \(AccA\), \(IOAccB\), \(IOAccA\) and \(S\) are opened at \(\zeta\) and \(\omega\zeta\). \(A\) and \(IOA\) are opened at \(\zeta\) and at \(\omega\zeta\) (F8, F31). Every other column is opened at \(\zeta\) only. Two opening points in total.

\textbf{Proof.} \texttt{CORRECTNESS\_PROOF.md} defines the clearing rule as a predicate \(\mathsf{Clear}\) and proves that the scalar checks and F1--F33 hold if and only if the receipt is \(\mathsf{Clear}\) of the book (soundness, completeness, uniqueness, and that the fills balance).

\textbf{What the constraints establish.} If F1--F33 hold on \(H\) and the scalar checks pass, then for the book \((B, A, IOB, IOA)\) underlying the commitments: \(V_{max}\) is the maximum executable volume; \([c,d]\) is the exact set of ticks attaining it; \(V_{min\Delta}\) is the minimum imbalance over \([c,d]\); \(p^*\) is the tick the tie-break rule of Section 7 selects, and \(s\) its long side; \(P\) is the long side\textquotesingle s volume at \(p^*\) and \(J\) the eligible short-side IO volume. The admitted IO volume, the residual, the traded volume and both fill ratios follow by public arithmetic. Conversely, an honest prover satisfies every constraint for any book whose volumes respect \(k_{in}\).

\begin{center}\rule{0.5\linewidth}{0.5pt}\end{center}

\subsection{\texorpdfstring{The No-Cross Case, \(V_{max} = 0\)}{The No-Cross Case, V\_\{max\} = 0}}\label{the-no-cross-case-v_max--0}

When no price crosses, \(Min\) is zero everywhere, the maximizing set is every tick, and the cliff equation is unsatisfiable. We treat this as a separate statement: the receipt is "no trade" and no \(c\), \(d\), \(p^*\) is disclosed.

One way to prove it is F1--F14 with \(V_{max} = 0\); F14 then forces \(Min \equiv 0\). We use a smaller, equivalent statement (revision 3, as implemented): commit only \(B\), \(A\), \(AccB\), \(AccA\), the limbs of \(B\) and \(A\) and the lookup columns for those limbs, and prove F1--F8, the lookup, and

\[AccA(X) \cdot AccB(X) = 0.\]

Both accumulators are sums of range-checked inputs, so they are non-negative integers below \(2^k\), and their product vanishes at a tick exactly when \(\min(AccA, AccB) = 0\) there. This needs \(4 + 2m_{in}\) value and limb columns in place of \(7 + 2m_{in} + 3m\), plus the lookup columns (\(M\), one helper per two looked-up values, \(U\), \(S\)): 17 columns for \(n = 8\), \(k_{in} = 8\). The IO columns take no part: IO orders cannot create a cross.

\begin{center}\rule{0.5\linewidth}{0.5pt}\end{center}

\subsection{Update Validity (revision 5, as implemented)}\label{update-validity-revision-5-as-implemented}

The clearing proof is about the book obtained by adding up every logged update. It does not show what each update did. Without a further check the exchange could log an update that subtracts a trader\textquotesingle s order after issuing the receipt (\texttt{SECURITY\_BOOK.md}, Section 6.1). The update-validity proof closes this. It is a second proof per batch, published with the clearing proof, with its own domain and columns.

\textbf{Statement.} The log holds steps \(\Delta_0, \ldots, \Delta_{m-1}\), each the change of one column commitment. For every \(j\) there are \(q_j \in [0, 2^{k_{in}})\) and a tick \(\ell_j\) such that \(\Delta_j\) commits to \(q_j L_{\ell_j}(X)\) up to a multiple of \(Z_H\) (or to a multiple of \(Z_H\) alone, an update that changes nothing).

\textbf{Steps per order.} An order logs four steps, one per column: its quantity on its own column and quantity zero on the other three (this hides the order\textquotesingle s side). So \(m\) is four times the number of orders, and a batch holds at most \(M/4\) orders.

\textbf{Domain.} \(K\), a multiplicative subgroup of size \(M \ge m\), with \(M = n\) or \(M \ge 2n\) a power of two, so that \(H \subseteq K\). Row \(j\) of \(K\) belongs to step \(j\). \(M\) is fixed at setup and is the largest number of orders a batch may hold. Write \(\kappa\) for the generator of \(K\), \(\sigma = M/n\), so that \(\kappa^{\sigma i} = \omega^i\).

\textbf{Protocol.}

\begin{enumerate}
\def\labelenumi{\arabic{enumi}.}
\item
  The exchange commits to \(Q\) with \(Q(\kappa^j) = q_j\) and \(T\) with \(T(\kappa^j) = \omega^{\ell_j}\) (rows \(j \ge m\): \(Q = 0\), \(T = 1\)), and to a fresh masked empty commitment \(\Delta_{mask}\) with the same proof of emptiness as a book opening (Section 12; \texttt{SECURITY\_BOOK.md} 3.1).
\item
  A challenge \(c\) is derived from the parameters, every \(\Delta_j\), \(\Delta_{mask}\), and the commitments to \(Q\) and \(T\).
\item
  Prover and verifier compute \(W = \Delta_{mask} + \sum_{j<m} c^j \Delta_j\). The prover knows the polynomial behind it: \(w(X) = \sum_i w_i L_i(X)\) plus the mask, with \(w_i = \sum_{j:\, \ell_j = i} c^j q_j\).
\item
  The proof system of Sections 2 and 8 is run on \(K\) with \(Q\), \(T\), \(W\) as columns whose commitments are already fixed, plus the limb columns of \(Q\) and the lookup columns.
\end{enumerate}

\textbf{Public polynomials on \(K\).} \(P(X)\) with \(P(\kappa^j) = c^j\) for \(j < m\) and 0 elsewhere; \(X_H(X)\) with value \(\omega^i\) at \(\kappa^{\sigma i}\) and 0 elsewhere; \(S_H(X)\), the indicator of the rows \(\kappa^{\sigma i}\); the range table \(T_w(X)\).

\textbf{Constraints, each on all of \(K\).}

\begin{longtable}[]{@{}lll@{}}
\toprule\noalign{}
\# & Name & Equation \\
\midrule\noalign{}
\endhead
\bottomrule\noalign{}
\endlastfoot
U1 & Quantity range & \(Q - \sum_l 2^{wl} \ell^Q_l = 0\), with every limb in the table (the lookup of Section 2) \\
U2 & Step side of the histogram & \(G(X)\,(\beta - T(X)) - P(X)\,Q(X) = 0\) \\
U3 & Tick side of the histogram & \(U_H(X)\,(\beta - X_H(X)) - S_H(X)\,W(X) = 0\) \\
U4 & Running sum & \(S_2(\kappa X) - S_2(X) - (G(X) - U_H(X)) = 0\), cyclic \\
\end{longtable}

\(G\), \(U_H\) and \(S_2\) are committed after \(\beta\), with the range lookup\textquotesingle s helpers. U2 to U4 prove \(\sum_j c^j q_j/(\beta - T_j) = \sum_{i<n} w_i/(\beta - \omega^i)\) in the same way as the lookup of Section 2.

\textbf{Why this is the statement.} Fix the log, \(Q\) and \(T\); then \(c\), then \(\beta\) are drawn.

\begin{itemize}
\item
  \emph{From U2 to U4.} As rational functions of \(\beta\), the two sides agree only if the residues agree at every point \(z\): \(\sum_{j:\,T_j = z} c^j q_j = w_i\) if \(z = \omega^i\) and \(0\) if \(z \notin H\). Otherwise the identity holds for at most \(M + n\) values of \(\beta\), plus at most \(M + n\) values at which a denominator vanishes and frees a helper (as in Section 2). No separate check that \(T_j \in H\) is needed: a step with \(T_j \notin H\) is forced to contribute nothing.
\item
  \emph{From the commitment.} \(W\) is computed by the verifier from the log. Its value on \(H\) is \(\sum_j c^j u_j\), where \(u_j\) is the value on \(H\) of the polynomial behind \(\Delta_j\) (\(\Delta_{mask}\) is zero on \(H\)). Binding of the commitment makes this the same \(w\) the proof opens.
\item
  \emph{Together.} \(\sum_j c^j \big(u_j - q_j e_{\ell_j}\big) = 0\) as vectors, where \(e_{\ell_j}\) is the unit vector at the tick of \(T_j\) (zero if \(T_j \notin H\)). The \(u_j\), \(q_j\), \(\ell_j\) were fixed before \(c\), so unless every \(u_j = q_j e_{\ell_j}\) this holds for at most \(M - 1\) values of \(c\).
\item
  \emph{U1} gives \(0 \le q_j < 2^{k_{in}}\).
\end{itemize}

So every logged step adds \(q_j \ge 0\) at one tick and nothing elsewhere, except with probability about \((3M + 2n + (2E_Q+2)M)/q\), where \(E_Q\) is the number of looked-up values of \(Q\). A step on a padding tick is not excluded here; the clearing proof\textquotesingle s padding constraints (F3, F4, F26, F27) reject the batch in that case.

\textbf{Hiding.} \(Q\) and \(T\) are fresh masked columns. \(W\) is opened at one point; the mask in \(\Delta_{mask}\) has four coefficients against two public functionals (its \(D\) and the opening), and its blinding four against two. Nothing else about the steps is opened.

\textbf{Cost.} Columns: \(Q\), \(T\), \(W\), \(\lceil k_{in}/w \rceil\) limbs, multiplicities, the range helpers, \(G\), \(U_H\), two running sums: 15 columns for \(k_{in} = 32\), \(w = 8\), so 17 commitments with the quotient and \(\Delta_{mask}\), and 1,700 bytes. Independent of the number of ticks and of the number of orders. The prover works on a domain of size \(M\); the verifier does one multi-scalar multiplication over the \(m\) logged steps. The setup must reach degree \(2M + 8\).

\begin{center}\rule{0.5\linewidth}{0.5pt}\end{center}

\subsection{Notes for Implementation}\label{notes-for-implementation}

\textbf{Zero-knowledge blinding.} The constraints above say nothing about hiding. A KZG10 commitment to a low-entropy column and its evaluation at \(\zeta\) both leak. For each committed polynomial \(P\) opened at \(m\) distinct points, the prover commits instead to \(P(X) + \rho(X) \cdot Z_H(X)\) with \(\rho\) uniformly random of degree \(m\) (that is, \(m+1\) random coefficients). The added term vanishes on \(H\), so every constraint is unaffected. Quotient degrees grow accordingly and the SRS must cover them.

\textbf{A cheaper strictness check.} F14 already proves \(V_{max} - Min \ge 0\) everywhere, so the cliffs only need \(V_{max} - Min \ne 0\). A single inverse column \(Inv(X)\) with \(Mask_C(X) \cdot \big[(V_{max} - Min(X)) \cdot Inv(X) - 1\big] = 0\) does this, and F19 gives the same shortcut for F21. This removes \(Slack\) and its \(k\) bit columns. The cost is a degree-\((3n-3)\) constraint. Under revision 3 that broke the uniform \(2n-2\) bound; since revision 4 the lookup already has degree \(3n-3\), so this is now a free optimization (it would remove one value column and \(m\) limb columns). It is not implemented; the slack form above is what the code proves.

\textbf{Inlining.} \(SurpB\), \(SurpA\), and \(\Delta\) are linear in other committed columns. Substituting them removes F9, F10, F18 and three commitments. We keep them committed because the per-trader check opens a surplus column directly.

\textbf{Why a logarithmic-derivative lookup and not Plookup.} The original file named Plookup (Gabizon and Williamson, ePrint 2020/315) for the input range check, with an identity that was not the Plookup identity. Revision 3 replaced it by bit columns. Revision 4 uses the lookup of Section 2 for these reasons:

\begin{enumerate}
\def\labelenumi{\arabic{enumi}.}
\item
  Plookup needs a committed sorted vector of length about \(2n\) per looked-up column and a grand product per column; logUp needs no sorting, and one multiplicity column and one running sum serve all 46 or so limb columns together.
\item
  Plookup checks \(n-1\) values over a size-\(n\) domain, so one tick needs separate treatment. The cyclic running sum checks all \(n\).
\item
  In either argument the table must fit in the domain. Limbs of \(w = \min(8, \log_2 n)\) bits remove the old consequence that the largest admissible volume per tick is about \(n\).
\end{enumerate}

\textbf{Bit columns versus limbs, measured.} At \(N = 1000\), \(k_{in} = 32\) (\(n = 1024\), \(k = 42\)), 8 threads, median of 5: bit columns give 352 commitments, a 22,924-byte proof, 319 ms to prove, 4.8 ms to verify; limbs with the lookup give 89 commitments, 6,124 bytes, 226 ms and 3.8 ms. The price is a second prover round, a quotient of degree about \(2n\) in place of \(n\), and an SRS twice as long.

\begin{center}\rule{0.5\linewidth}{0.5pt}\end{center}

\subsection{Map from the Original List}\label{map-from-the-original-list}

\begin{longtable}[]{@{}lll@{}}
\toprule\noalign{}
Original & Status & Now \\
\midrule\noalign{}
\endhead
\bottomrule\noalign{}
\endlastfoot
\#1, \#2 Plookup range & Replaced by limb decomposition with a logUp lookup (Section 2); stated identity was not Plookup & F1, F2, F23 \\
--- & Added: padding ticks carry no volume & F3, F4 \\
\#3, \#4 Init & Kept, written with \(L_{n-1}\), \(L_0\) & F5, F6 \\
\#5 Demand transition & Kept & F7 \\
\#6 Supply transition & \textbf{Fixed:} \(A(X) \to A(\omega X)\) & F8 \\
--- & \textbf{Added:} surplus definitions were missing & F9, F10 \\
\#7 Min mutual exclusivity & Kept, rewritten through the surpluses & F11 \\
\#8 Ceiling & Kept; width is the fixed \(k\) & F14 \\
\#9, \#10 \(Mask_P\) booleanity, position & Removed: the mask is public, nothing to prove & --- \\
\#11, \#12 Endpoints & Kept & F15, F16 \\
\#13, \#14 \(InMCV\) & Removed: \(InMCV\) was never tied to \(Min\) and constrains nothing & --- \\
\#15, \#16 Surplus non-negativity & Kept & F12, F13 \\
\#17 \(\Delta\) definition & Kept & F18 \\
\#18--\#20 \(ChkD\) & Removed: with public \(p^*\), their only effect is \(p^* \in [c,d]\), now a scalar check & Section 1 \\
\#21, \#22 \(Mask_V\) & Removed, same reason & --- \\
\#23 Delta floor & Kept & F19 \\
\#24 Valley pin & Kept & F20 \\
--- & \textbf{Added:} tie-break among equal-imbalance ticks & F21 \\
\#25, \#26 \(SD\) & Removed: implied by \#24 & --- \\
\#27, \#28 \(Mask_C\) booleanity, containment & Removed: the mask is public & --- \\
\#29, \#30 Cliffs & Merged; boundary cases fixed; gated by \(Mask_C\) & F17 \\
\#31, \#32 Slack non-negativity & Merged into one column & F22 \\
\#33 Booleanity & Replaced by table membership of the limbs & F23 \\
Cliff value openings & Removed: they leak depth & --- \\
\end{longtable}

Phases 1 and 2 go from 33 constraints to 24 identities (F1--F23, with F20 split in two). Phase 3 adds ten (F24--F33), for 34 in total. Every sentence in the draft that says "33 constraints" depends on this.

\begin{center}\rule{0.5\linewidth}{0.5pt}\end{center}

\subsection{Outside the Proof}\label{outside-the-proof}

\begin{itemize}
\item
  \textbf{Per-order rounding.} The proof fixes the aggregate quantities and the two fill ratios. Fills are integers, so \(q \cdot (P-R)/P\) must be rounded, and distributing the leftover shares depends on the other orders at the tick. We do not prove a rounding rule. The venue must state one; the simplest is to round every fill down and leave the remainder unexecuted.
\item
  \textbf{Binding to the published book.} The constraints are stated over \(\mathsf{cm}[B]\), \(\mathsf{cm}[A]\), \(\mathsf{cm}[IOB]\), \(\mathsf{cm}[IOA]\). They prove a statement about the traders\textquotesingle{} orders only if those are the commitments the update receipts built up during the batch. This is the trader-side check, not part of the clearing proof.
\item
  \textbf{Orders the exchange inserts itself.} Nothing distinguishes an exchange-authored order from a trader\textquotesingle s. This is true of any venue.
\item
  \textbf{Orders the exchange removes: now inside the proofs.} The clearing proof alone does not show that a logged update adds volume. The update-validity proof of Section 9b does (added 2026-11-18, after \texttt{SECURITY\_BOOK.md} Section 6.1 found the gap). A batch is valid only if both proofs verify.
\item
  \textbf{Order side: hidden since 2026-11-18.} Every order logs a step on each of the four columns, three of them with quantity zero, so the log does not show an order\textquotesingle s side or type (\texttt{SECURITY\_BOOK.md} Section 6.2). The number and timing of orders remain public.
\item
  \textbf{One proof per book.} The off-domain mask of a book column covers one batch-close proof. The implementation derives the prover\textquotesingle s randomness from the private state, so closing the same batch again reproduces the same proof, and it refuses orders after a close (\texttt{SECURITY\_BOOK.md} Section 6.3). The rule itself remains an operating requirement on the exchange.
\end{itemize}

\begin{center}\rule{0.5\linewidth}{0.5pt}\end{center}

\subsection{Design Decisions}\label{design-decisions}

The author delegated these four. Each is settled as follows.

\begin{enumerate}
\def\labelenumi{\arabic{enumi}.}
\item
  \textbf{Tie-break rule.} Not "lowest tick". We use the rule Phase 3 forces: the lowest minimizing tick whose at-price long-side volume covers the imbalance (Section 7). It is the only rule under which "orders strictly better than \(p^*\) execute in full" is always satisfiable.
\item
  \textbf{Range checks by limbs and one lookup.} Every inequality is a range check, and every range check is a limb decomposition whose limbs go through one shared logUp lookup (Section 2). Plookup is dropped; so are the bit columns of revision 3. The protocol has one inequality gadget, which the paper can present as "range check = limbs + lookup".
\item
  \textbf{Hidden cliffs.} The executable volume at \(c-1\) and \(d+1\) stays private; strictness is proved with the slack column. The receipt should disclose only what the auction outcome requires.
\item
  \textbf{No-cross case.} A batch with \(V_{max} = 0\) is a separate, smaller statement (Section 9).
\end{enumerate}
